\documentclass[10pt,journal]{IEEEtran}

\usepackage{iftex}
\ifPDFTeX
  \usepackage[utf8]{inputenc}
  \usepackage[T1]{fontenc}
\fi
\usepackage{amsmath,amssymb,mathtools,bm}
\usepackage{newtxtext,newtxmath}
\usepackage{textcomp}
\usepackage{booktabs,array}
\usepackage{graphicx}
\usepackage{pgfplots}
\usepackage{cite}
\usepackage[final]{microtype}
\usepackage[breaklinks=true,hidelinks]{hyperref}
\interdisplaylinepenalty=2500
\emergencystretch=2em
\allowdisplaybreaks
\pgfplotsset{compat=1.18}
\definecolor{oublue}{HTML}{1F77B4}
\definecolor{shoorange}{HTML}{FF7F0E}
\definecolor{colgreen}{HTML}{2CA02C}
\definecolor{carmared}{HTML}{D62728}
\newcommand{\E}{\mathbb E}
\newcommand{\R}{\mathbb R}
\newcommand{\T}{^{\mathsf T}}
\newcommand{\PSD}{S_x}
\newcommand{\Hz}{\,\mathrm{Hz}}

\title{Other Stochastic Models for Marine Motion in Kalman Filters}

\author{Mikhail S. Grushinskiy%
\thanks{Independent researcher, Bareboat Necessities project, Fair Lawn, NJ, USA.
Correspondence: mgrouch@users.github.com.}}

\begin{document}
\maketitle

\begin{abstract}
A low-cost marine inertial measurement unit senses specific force, not
displacement. Twice integrating corrected acceleration is vulnerable to
low-frequency noise, attitude error, and accelerometer bias. In a broad-banded
sea state with one dominant encounter frequency, a stationary stochastic
displacement prior may give an appropriate alternative to unrestrained
integration. This research note develops a hierarchy of \emph{linear}
continuous-time stochastic differential equations: an ordinary
Ornstein--Uhlenbeck (OU) displacement, a white-noise-forced damped oscillator,
an oscillator driven by colored OU forcing, and spectrally shaped continuous-time
autoregressive moving-average (CARMA) processes. We derive power spectra,
acceleration regularity requirements, and exact Kalman state-space
realizations. Two reproducible comparison figures illustrate time traces and
spectral shapes. Particular attention is paid to CARMA(3,0), (4,1), and (5,2),
which permit mean-square acceleration, in contrast to CARMA(2,0) and (4,2).
These are \emph{candidate models}, not implemented Ocean-IMU methods or claims
of improved heave accuracy. The note ends with a controlled experimental
program for model selection, adaptation, and possible integration.
\end{abstract}

\begin{IEEEkeywords}
marine heave, Kalman filtering, continuous-time stochastic processes,
CARMA, Ornstein--Uhlenbeck process, spectral factorization, inertial navigation
\end{IEEEkeywords}

\section{The Reconstruction Problem}
Let \(x(t)\) denote the vessel's vertical displacement relative to a
local equilibrium (not absolute Earth-fixed height), with velocity \(v=\dot x\)
and acceleration \(a=\ddot x\). The idealized measurement made after rotating
specific force into a level frame and compensating gravity is
\begin{equation}
y_a(t)=a(t)+b_a(t)+n_a(t)+\epsilon_R(t),
\label{eq:measurement}
\end{equation}
where \(b_a\) denotes residual accelerometer bias, \(n_a\) finite-bandwidth
measurement noise, and \(\epsilon_R\) projection error due to imperfect
attitude and other unmodeled dynamics. The actual body-frame relationship is
\begin{equation}
\mathbf f_m^B=R_{WB}^{\mathsf T}
  (\mathbf a^W-\mathbf g^W)+\mathbf b_a^B+\mathbf n_a^B.
\label{eq:body-meas}
\end{equation}
Thus (\ref{eq:measurement}) is a useful \emph{conditional} vertical subproblem,
not a substitute for the coupled attitude--translation measurement model.

We seek a stationary displacement process with a broad power spectrum but
one prominent frequency \(f_d>0\), changing amplitudes and phases, and
sensible suppression of poorly observed very-low-frequency modes.
For a vessel, the appropriate target is the \emph{vessel heave response}
spectrum. The incident JONSWAP/PM elevation spectrum is not generally identical
to it: encounter kinematics, directional spread, and the vessel response
amplitude operator modify the heave spectrum.

\section{A Common Linear SDE Description}
All candidates admit an augmented Markov state \(z\in\R^p\):
\begin{equation}
dz=A z\,dt+G\,dW_t,\qquad x=Cz,
\label{eq:linear-sde}
\end{equation}
where \(A\) and \(G\) are constant during one identified sea-state interval
and \(W_t\) is standard Brownian motion. A Hurwitz \(A\), stationary initial
state, and Gaussian driving noise define a stationary Gaussian output.
Its state covariance satisfies
\begin{equation}
A P_\infty+P_\infty A\T+GG\T=0.
\label{eq:lyapunov}
\end{equation}
The two-sided displacement spectral density, with
\(\mathrm{Var}(x)=(2\pi)^{-1}\int_{\R}\PSD(\omega)\,d\omega\), is
\begin{equation}
\PSD(\omega)=
C(i\omega I-A)^{-1}GG\T
(-i\omega I-A\T)^{-1}C\T.
\label{eq:state-psd}
\end{equation}
A finite-dimensional \emph{state} can be Markov even when the scalar
displacement output is not.

\section{The First Three Models}
\subsection{OU Displacement: CARMA(1,0)}
The ordinary OU process,
\begin{equation}
dx=-\lambda x\,dt+\sigma\,dW_t,\quad
\PSD(\omega)=\frac{\sigma^2}{\lambda^2+\omega^2},
\label{eq:ou}
\end{equation}
has maximum power at zero frequency; it has no resonant band or
mean-square derivative. It is a useful model for a slowly correlated
\emph{disturbance}, but not a good standalone prior for oscillatory
heave displacement.

\subsection{White-Noise-Driven Oscillator: CARMA(2,0)}
For natural frequency \(\omega_0\) and damping ratio \(\zeta>0\),
\begin{align}
dx &= v\,dt,\nonumber\\
dv &=(-\omega_0^2x-2\zeta\omega_0v)\,dt+\sigma\,dW_t,
\label{eq:sho}
\end{align}
with spectrum
\begin{equation}
\PSD(\omega)=
\frac{\sigma^2}
{(\omega_0^2-\omega^2)^2+
 (2\zeta\omega_0\omega)^2}.
\label{eq:sho-psd}
\end{equation}
For \(0<\zeta<1/\sqrt2\) its nonzero-frequency peak is at
\(\omega_{\rm peak}=\omega_0\sqrt{1-2\zeta^2}\).
It produces random-amplitude and random-phase oscillations; however,
the formal acceleration has a white-noise component. Velocity is
mean-square differentiable only if the idealized forcing is additionally
smoothed. A real accelerometer is finite-bandwidth, but this distinction
matters when acceleration is itself the continuous-time observation.

\subsection{Colored-Forcing Oscillator: CARMA(3,0)}
Drive the same oscillator through an OU force \(u\):
\begin{equation}
\begin{aligned}
dx&=v\,dt,\\
dv&=(-\omega_0^2x-2\zeta\omega_0v+u)\,dt,\\
du&=-\lambda u\,dt+\sigma_u\,dW_t .
\end{aligned}
\label{eq:colored}
\end{equation}
For \(z=[x,v,u]\T\), this is (\ref{eq:linear-sde}) with
\begin{equation}
A=\begin{bmatrix}
0&1&0\\
-\omega_0^2&-2\zeta\omega_0&1\\
0&0&-\lambda
\end{bmatrix},\quad
G=\begin{bmatrix}0\\0\\\sigma_u\end{bmatrix}.
\label{eq:colored-matrix}
\end{equation}
Its acceleration,
\(a=-\omega_0^2x-2\zeta\omega_0v+u\), is a
finite-variance, continuous stochastic process. Eliminating \(u\)
gives a third-order linear stochastic operator:
\begin{equation}
(D+\lambda)(D^2+2\zeta\omega_0D+\omega_0^2)x
=\sigma_u\,\xi,\quad D=d/dt,
\label{eq:colored-carma}
\end{equation}
where \(\xi\) is formal Gaussian white noise. The
spectrum is (\ref{eq:sho-psd}) multiplied by
\((\lambda^2+\omega^2)^{-1}\) with the appropriate driving
intensity. Relative to the white-forced oscillator, it attenuates
high-frequency displacement and allows an ordinary linear
acceleration measurement map.

\section{The CARMA Generalization}
A Gaussian CARMA(\(p,q\)) model \cite{brockwell01,brockwell04} has
transfer function \(H(s)=\sigma b(s)/a(s)\), where
\begin{align}
a(s)&=s^p+a_1s^{p-1}+\cdots+a_p,\nonumber\\
b(s)&=b_0+b_1s+\cdots+b_qs^q,\qquad 0\le q<p,
\label{eq:carma-poly}
\end{align}
with no cancellation of leading polynomial degree. In symbolic
stochastic-operator notation,
\begin{equation}
a(D)x=\sigma b(D)\xi,\qquad
\PSD(\omega)=\sigma^2\frac{|b(i\omega)|^2}{|a(i\omega)|^2}.
\label{eq:carma-psd}
\end{equation}
All roots of \(a\) must have negative real parts for a causal,
stationary state solution. Its poles set damping and resonances;
the numerator provides frequency-dependent shaping. A real-rational
nonnegative target spectrum can be approximated by a stable
spectral factor over a chosen band, but a fixed finite-order
CARMA model cannot reproduce an arbitrary spectrum exactly.
CARMA descriptions of random seas have precedents, including
JONSWAP spectral approximations \cite{dostal11}; those
examples do not establish accuracy for vessel heave inferred
from an IMU.

\subsection{The Acceleration-Regularity Test}
If the numerator has degree \(q\) and the denominator degree \(p\),
\begin{equation}
\PSD(\omega)=O(|\omega|^{-2(p-q)}),
\qquad |\omega|\rightarrow\infty.
\label{eq:tail}
\end{equation}
A mean-square second derivative requires
\(\int_{\R}\omega^4\PSD(\omega)d\omega<\infty\).
For ordinary finite-degree CARMA spectra without leading-degree
cancellation, this is equivalent to
\begin{equation}
\boxed{p-q\ge 3.}
\label{eq:acc-criterion}
\end{equation}
It is a \emph{continuous-time, ideal acceleration} criterion.
A finite-bandwidth accelerometer can still observe a filtered
CARMA(2,0) or CARMA(4,2) signal, provided the measurement filtering
is explicitly modeled.

\begin{table}[t]
\centering
\caption{CARMA displacement candidates: finite-variance ideal
acceleration means \(\E|\ddot x|^2<\infty\).}
\label{tab:regularity}
\footnotesize
\setlength{\tabcolsep}{5pt}
\begin{tabular}{@{}lcc@{}}
\toprule
Displacement prior & High-frequency \(\PSD\) & Accel.\ exists?\\
\midrule
CARMA(1,0) -- ordinary OU & \(\omega^{-2}\) & \textbf{No}\\
CARMA(2,0) -- white oscillator & \(\omega^{-4}\) & \textbf{No}\\
CARMA(3,0) -- colored oscillator & \(\omega^{-6}\) & \textbf{Yes}\\
CARMA(4,2) -- spectral example & \(\omega^{-4}\) & \textbf{No}\\
CARMA(4,1) -- candidate & \(\omega^{-6}\) & \textbf{Yes}\\
CARMA(5,2) -- candidate & \(\omega^{-6}\) & \textbf{Yes}\\
\bottomrule
\end{tabular}
\end{table}

\subsection{Fourth- and Fifth-Order Candidates}
Consider two stable damped-oscillator factors
\begin{equation}
L_j(s)=s^2+2\zeta_j\omega_js+\omega_j^2,\qquad j=1,2.
\end{equation}
A four-state acceleration-compatible candidate is
\begin{equation}
H_{4,1}(s)=\frac{\sigma\kappa s}{L_1(s)L_2(s)}.
\label{eq:41}
\end{equation}
This has \(\PSD\propto\omega^2\) near zero and
\(\PSD\propto\omega^{-6}\) at high frequency. The first
pole pair can govern the dominant encounter resonance; the second
changes the shoulder and bandwidth.

Adding another stable real pole while raising the numerator degree
gives the five-state proposal
\begin{equation}
H_{5,2}(s)=
\frac{\sigma\kappa s^2}{(s+\lambda)L_1(s)L_2(s)},
\label{eq:52}
\end{equation}
with \(\PSD\propto\omega^4\) at low frequency and
\(\PSD\propto\omega^{-6}\) at high frequency.
This stronger low-frequency rejection is an \emph{assumption},
not a free improvement: it risks filtering out true long-period
heave.

For comparison, the plotted four-state CARMA(4,2) example is
\begin{equation}
H_{4,2}(s)=\frac{\sigma\kappa s^2}{L_1(s)L_2(s)},
\label{eq:42}
\end{equation}
which can fit a broad high-frequency shoulder but has
\(\PSD=O(\omega^{-4})\); it therefore fails
(\ref{eq:acc-criterion}). Its inclusion in the figures is
illustrative, not an endorsement as the ideal-acceleration model.

\section{Illustrative Time Traces and Spectra}
Figures~\ref{fig:time} and \ref{fig:spectra} reproduce the
four-model comparisons motivating this note. They use illustrative
model parameters, not identified vessel data:
\(f_0=0.29\Hz\), \(\zeta=0.22\), OU displacement
\(\lambda=2\pi(0.12\Hz)\), forcing OU
\(\lambda_u=2\pi(0.40\Hz)\). For CARMA(4,2), the
oscillator pairs are \((f_1,\zeta_1)=(0.27\Hz,0.20)\)
and \((f_2,\zeta_2)=(0.56\Hz,0.60)\).
Original linear-SDE simulations were exactly discretized
using matrix exponentials and stationary Gaussian initial conditions.
The plot table is sampled at \(5\Hz\) from a \(25\Hz\)
realization, with each displayed signal divided by its
theoretical stationary standard deviation. The spectral
figure uses analytic spectra, each normalized to its own peak;
the values below \(-70\) dB are clipped for display.
A single realization does not establish estimation performance.

\begin{figure*}[t]
\centering
\begin{tikzpicture}
\begin{axis}[
 width=.90\textwidth,height=6.0cm,
 xmin=0,xmax=65,ymin=-3.0,ymax=12.8,
 xlabel={Time (s)},
 ylabel={Normalized displacement (stacked)},
 ytick={0,3,6,9},
 yticklabels={CARMA 4/2,Colored force,White force,Ordinary OU},
 y tick label style={font=\scriptsize},
 tick label style={font=\scriptsize},
 grid=major,grid style={gray!17},
 axis x line*=bottom,axis y line*=left,
 every axis plot/.append style={thin,no marks},
]
\addplot[oublue] table [col sep=comma,x=t,y=ou]
  {carma-heave-time.csv};
\addplot[shoorange] table [col sep=comma,x=t,y=sho]
  {carma-heave-time.csv};
\addplot[colgreen] table [col sep=comma,x=t,y=colored]
  {carma-heave-time.csv};
\addplot[carmared] table [col sep=comma,x=t,y=carma42]
  {carma-heave-time.csv};
\end{axis}
\end{tikzpicture}
\caption{Representative displacement traces, stacked at offsets
9, 6, 3, and 0 (in dimensionless RMS units) respectively.
The OU process is irregular but non-oscillatory; oscillator and
CARMA realizations fluctuate in amplitude and phase around
a dominant period. The four traces are \emph{different} realizations,
not the result of a shared physical forcing.}
\label{fig:time}
\end{figure*}

\begin{figure*}[t]
\centering
\begin{tikzpicture}
\begin{axis}[
 width=.90\textwidth,height=6.0cm,
 xmin=0,xmax=1.3,ymin=-70,ymax=4,
 xlabel={Frequency (Hz)}, ylabel={Displacement PSD (dB re.\ own peak)},
 xtick={0,0.2,...,1.2},
 ytick={-70,-60,...,0},
 tick label style={font=\scriptsize},
 legend style={font=\scriptsize,at={(.98,.98)},anchor=north east,
   fill=white,draw=gray!40},
 grid=major,grid style={gray!17},
 every axis plot/.append style={thick,no marks},
]
\addplot[oublue] table [col sep=comma,x=f,y=ou]
  {carma-heave-spectra.csv};
\addlegendentry{Ordinary OU}
\addplot[shoorange] table [col sep=comma,x=f,y=sho]
  {carma-heave-spectra.csv};
\addlegendentry{White-forced oscillator}
\addplot[colgreen] table [col sep=comma,x=f,y=colored]
  {carma-heave-spectra.csv};
\addlegendentry{Colored-forced oscillator}
\addplot[carmared] table [col sep=comma,x=f,y=carma42]
  {carma-heave-spectra.csv};
\addlegendentry{CARMA(4,2)}
\end{axis}
\end{tikzpicture}
\caption{Analytic continuous-time spectral shapes for the same
four models. A resonant peak around \(0.28\Hz\) is absent
from the OU displacement. Colored forcing gives a steeper
high-frequency tail. The CARMA(4,2) red curve displays
strong low-frequency suppression, but has only an
\(\omega^{-4}\) asymptotic tail and thus no finite-variance
ideal acceleration. Normalization removes absolute power information.}
\label{fig:spectra}
\end{figure*}

\section{Exact Kalman-Compatible State-Space Form}
For the polynomial convention in (\ref{eq:carma-poly}), define
the controllable companion realization
\begin{equation}
A_c=\begin{bmatrix}
0&1&0&\cdots&0\\
0&0&1&\cdots&0\\
\vdots&&&\ddots&\vdots\\
0&0&0&\cdots&1\\
-a_p&-a_{p-1}&-a_{p-2}&\cdots&-a_1
\end{bmatrix},
\label{eq:companion}
\end{equation}
\begin{equation}
G_c=\sigma e_p,\quad
C_c=\begin{bmatrix}b_0&b_1&\cdots&b_q&0&\cdots&0\end{bmatrix}.
\label{eq:companion-out}
\end{equation}
Then \(d z=A_cz\,dt+G_c\,dW_t\), \(x=C_cz\)
is a linear It\^o SDE with the intended transfer function.
For \(p-q\ge3\), \(C_cG_c=C_cA_cG_c=0\) and
\begin{equation}
\boxed{x=C_cz,\qquad v=C_cA_cz,\qquad
a=C_cA_c^2z.}
\label{eq:xva}
\end{equation}
Those are state outputs, not numerical derivatives of noisy
displacement samples.

At sampling interval \(\Delta t\), the exact discrete
model is
\begin{align}
z_{k+1}&=Fz_k+w_k,\qquad F=e^{A_c\Delta t},\nonumber\\
Q&=\E[w_kw_k\T]
 =\int_0^{\Delta t}e^{A_c\tau}G_cG_c\T
                  e^{A_c\T\tau}\,d\tau.
\label{eq:exact-discrete}
\end{align}
For stationary coefficients and covariance (\ref{eq:lyapunov}),
\begin{equation}
Q=P_\infty-FP_\infty F\T
\label{eq:stationary-Q}
\end{equation}
is an equivalent exact formula. Matrix-exponential/Van Loan
discretization avoids Euler sampling errors
\cite{sarkka23}.

A reduced vertical filter may augment \(z\) with a bias \(b\):
\begin{align}
\begin{bmatrix}z\\b\end{bmatrix}_{k+1}
&=\begin{bmatrix}F&0\\0&F_b\end{bmatrix}
  \begin{bmatrix}z\\b\end{bmatrix}_k+w_k^+,\nonumber\\
y_{a,k}&=\begin{bmatrix}C_cA_c^2&1\end{bmatrix}
\begin{bmatrix}z\\b\end{bmatrix}_k+n_k^+ .
\label{eq:kalman-y}
\end{align}
The choice \(F_b=1\) denotes a random-walk residual bias;
other bias priors can be examined without claiming they match
the deployed bias law. The \(n_k^+\) covariance must account
for projected accelerometer noise, attitude uncertainty,
anti-aliasing, and any process/measurement correlation caused
by sampling or preprocessing.

In a coupled quaternion MEKF, the vertical force residual
also depends on attitude error and body-frame bias.
Its Jacobian must retain those cross terms; inserting
(\ref{eq:kalman-y}) unchanged into a full nonlinear
attitude filter would be incorrect. Conditional on a fixed
or externally supplied orientation and fixed CARMA
parameters, however, the translational subproblem is linear
Gaussian and has the ordinary Kalman recursion.

\subsection{Bias, Identifiability, and Adaptation}
Acceleration has frequency response
\(A(\omega)=-\omega^2 X(\omega)\). Near zero frequency it
therefore carries little displacement information, and a
constant or slowly varying accelerometer bias can be
confounded with low-frequency motion. A stationary prior
regularizes the estimate but cannot manufacture missing
measurement information. Strong zero-frequency suppression
can improve apparent drift while hiding real long-period
heave. Innovation whiteness and small estimated covariance
alone do not prove accuracy.

Stable coefficients \(\theta=(\omega_j,\zeta_j,\lambda,
b_j,\sigma)\) may be estimated from a windowed spectral fit
or a predictive innovations likelihood, then held fixed
over short intervals. The \emph{conditional} state estimator
remains linear between coefficient changes; jointly
estimating random or state-dependent frequencies is a
nonlinear estimation problem. A parameter-switch policy
needs explicit covariance continuity, initial-condition
handling, and validation under changing seas; blindly
restarting a stationary filter at each frequency update
would erase information.

\section{Research Program for Ocean-IMU}
The following constitutes a \emph{proposed} study, not
experimental results or a deployment change.

\subsection{Model Identification}
Estimate target \emph{vessel heave} power spectra using
truth displacement when available, or a carefully
characterized external reference. Fit CARMA(3,0), (4,1),
and (5,2) over a relevant encounter-frequency band with
stable pole constraints and nonnegative driving intensity.
A log-spectral residual can weight weakly energized bands,
but fitting only spectral shape can obscure physical
amplitude; validate both PSD shape and integrated variance.
Compare held-out likelihood, bias/variance, pole sensitivity,
and parameter identifiability. CARMA(2,0) and (4,2)
remain useful finite-bandwidth baselines with explicit
sensor filtering, not ideal-acceleration candidates.

\subsection{Matched Heave Reconstruction Ablation}
Use the \emph{same} synthetic or recorded IMU force,
attitude history, temperature, bias history, anti-alias
filter, timestamps, and ground-truth heave for every
candidate. Include quiet water; narrow and broad JONSWAP,
PM and nonlinear/Stokes-like responses; varying encounter
frequency; crossing seas; long-period swell; and
high-frequency vibration. Compare a controlled OU-driven
baseline, CARMA(3,0), (4,1), (5,2), and the deployed
OU--III estimator \emph{as actually implemented}, with no
silent changes to its gates, pseudo-measurements, tuning,
or magnetic-service behavior.

Quantify displacement RMS and bias, normalized heave
error where significant wave height is defined,
frequency-dependent gain/phase, low-frequency spectral
leakage, innovation consistency, false heave at rest,
transient recovery, numerical condition, and embedded
memory/runtime. Include held-out vessel responses rather
than merely simulating data from the same CARMA model
being evaluated. A controlled stress test must assess
whether more aggressive low-frequency regularization
simply suppresses true long-period displacement.

\subsection{Relation to the Shipping OU--III Stability Study}
The statistical displacement prior in this proposal is
\emph{not a law for the actual vessel trajectory}.
In particular, a nondegenerate stationary Gaussian
CARMA realization is not almost surely globally bounded
over an infinite horizon. Even when \(p-q=3\) gives
finite-variance continuous acceleration, its sample-path
jerk is not subject to the shipping proof's deterministic
bounded-jerk premise. A stochastic simulation does not
establish admissibility, finite capture, or regional
practical stability of OU--III. Future replacements would
require their own process-model mismatch, observability,
covariance, and full estimator stability analyses.
Nothing in this article revises the shipping estimator
or its proof assumptions.

\section{Conclusion}
A plain OU displacement prior misses resonant marine
heave; a white-forced second-order oscillator represents
the resonance but not finite-variance ideal acceleration.
Colored forcing yields a CARMA(3,0) baseline, while
CARMA(4,1) and (5,2) allow additional low-/high-frequency
spectral shaping with ordinary acceleration outputs and
exact discrete-time Kalman transitions. CARMA(4,2)
illustrates that greater spectral flexibility alone
does not guarantee physically appropriate derivatives.
The next step is not to declare a new best filter, but
to test whether spectral fit improvements survive
bias-contaminated, vessel-response-aware, held-out
heave reconstruction at realistic IMU sampling rates.

\section*{Acknowledgment and Status}
This is a supervised AI-assisted research proposal
for future evaluation. The author retains responsibility
for mathematical and experimental verification.
No CARMA-based Ocean-IMU runtime results are claimed.

\begin{thebibliography}{9}
\bibitem{brockwell01}
P. J. Brockwell, ``Continuous-time ARMA processes,''
in \emph{Handbook of Statistics}, vol. 19,
Elsevier, 2001, pp. 249--276,
doi: 10.1016/S0169-7161(01)19011-5.
\bibitem{brockwell04}
P. J. Brockwell, ``Representations of continuous-time
ARMA processes,'' \emph{J. Appl. Probab.},
vol. 41, no. A, pp. 375--382, 2004,
doi: 10.1239/jap/1082552212.
\bibitem{dostal11}
L. Dostal and E. Kreuzer,
``Probabilistic approach to large amplitude ship rolling in
random seas,'' \emph{Proc. IMechE, Part C: J. Mechanical
Engineering Science}, vol. 225, no. 10,
pp. 2464--2476, 2011,
doi: 10.1177/0954406211414523.
\bibitem{sarkka23}
S. S\"arkk\"a and L. Svensson,
\emph{Bayesian Filtering and Smoothing}, 2nd ed.,
Cambridge University Press, 2023,
doi: 10.1017/9781108917407.
\end{thebibliography}
\end{document}